\subsection{Равномерное движение}
%%%%%%%%%%%%%%%%%%%%%%%%%%%%%%%%%%%%%%%%%%%%%%%%%%%%
%%%%%%%%%%%%%%%%%%%%%%%%%%%%%%%%%%%%%%%%%%%%%%%%%%%%
%%%%%%%%%%%%%%%%%%%%%%%%%%%%%%%%%%%%%%%%%%%%%%%%%%%%
\z{Какое время двигался автомобиль, если на первую треть своего пути он затратил на $10\,\text{мин}$ больше, чем на оставшийся путь, скорость на котором была втрое больше, чем на первой трети пути?}
%
{$t = 50\,\text{мин}$}

%%%%%%%%%%%%%%%%%%%%%%%%%%%%%%%%%%%%%%%%%%%%%%%%%%%%
\z{Автобус движется к мосту со скоростью $72\,\text{км/ч}$. В начальный момент расстояние от него до моста равно $200\,\text{м}$. На каком расстоянии от моста будет автобус через $50\,\text{с}$, если длина моста $400\,\text{м}$?}
%
{$S = 400\,\text{м}$}

%%%%%%%%%%%%%%%%%%%%%%%%%%%%%%%%%%%%%%%%%%%%%%%%%%%%
\z{Товарный поезд движется со скоростью $v_1 = 36\,\text{км/ч}$. Спустя время $\tau = 30\,\text{мин}$ с той же станции в том же направлении вышел экспресс со скоростью $v_2 = 72\,\text{км/ч}$. Через время $t$ после выхода товарного поезда и на каком расстоянии $s$ от станции экспресс нагонит товарный поезд? Задачу решить, используя закон движения.}
%
{$t = 1\,\text{ч}; \quad s = 36\,\text{км}$}

%%%%%%%%%%%%%%%%%%%%%%%%%%%%%%%%%%%%%%%%%%%%%%%%%%%%
\z{Из пункта $A$ выехал велосипедист со скоростью $v_1 = 25\,\text{км/ч}$. Спустя время $t_0 = 6\,\text{мин}$ из пункта $B$, находящегося на расстоянии $L = 10\,\text{км}$ от пункта $A$, навстречу велосипедисту вышел пешеход. За время $t_2 = 50\,\text{с}$ пешеход прошёл такой же путь, какой велосипедист проехал за $t_2 = 10\,\text{с}$. На каком расстоянии $s$ от пункта $A$ встретятся пешеход и велосипедист?}
%
{$s = 8.75\,\text{км}$}

%%%%%%%%%%%%%%%%%%%%%%%%%%%%%%%%%%%%%%%%%%%%%%%%%%%%
\z{Из пунктов $A$ и $B$ одновременно навстречу друг другу выехали две машины. Через некоторое время они встретились и продолжили своё движение. Первая машина пришла в пункт назначения через $t_1 = 4\,\text{ч}$ после встречи, вторая — через $t_2 = 1\,\text{ч}$. Через какое время после выхода из пунктов $A$ и $B$ машины встретились?}
%
{$t = \sqrt{t_1 \cdot t_2} = 2\,\text{ч}$}

%%%%%%%%%%%%%%%%%%%%%%%%%%%%%%%%%%%%%%%%%%%%%%%%%%%%
\z{На прямой дороге находятся велосипедист, мотоциклист и пешеход между ними. В начальный момент времени расстояние от пешехода до велосипедиста в $2$ раза меньше, чем до мотоциклиста. Велосипедист и мотоциклист начинают двигаться навстречу друг другу со скоростями $20\,\text{км/ч}$ и $60\,\text{км/ч}$ соответственно. В какую сторону и с какой скоростью должен идти пешеход, чтобы встретиться с велосипедистом и мотоциклистом в месте их встречи?}
%
{$v = 6.7\,\text{км/ч}; \quad \text{к велосипедисту}$.  }

%%%%%%%%%%%%%%%%%%%%%%%%%%%%%%%%%%%%%%%%%%%%%%%%%%%%
\z{Экспериментатор Глюк наблюдал с безопасного расстояния за движением грозовой тучи. Увидев первую молнию, он засёк время и обнаружил, что услышал гром от неё только через $t_1 = 20\,\text{с}$. Через $\tau_1 = 3\,\text{мин}$ после первой вспышки произошла вторая, а гром грянул с опозданием на $t_2 = 5\,\text{с}$. Подождав ещё $\tau_2 = 4\,\text{мин}$ после второй вспышки, Глюк увидел, как сверкнула последняя молния, и услышал звук грома от неё через $t_3 = 20\,\text{с}$. Определите скорость тучи $v$ и минимальное расстояние $h$ от Глюка за время наблюдения. Скорость звука в воздухе $u \approx 330\,\text{м/с}$, скорость света $c = 3 \cdot 10^8\,\text{м/с}$.}
%
{$v = 32\,\text{м/с}; \quad h = 1.4\,\text{км}$. }